\documentclass[10pt]{article}
\usepackage[T1]{fontenc}
\usepackage[utf8]{inputenc}
\usepackage[french]{babel}
\usepackage{newpxtext,newpxmath}
\usepackage[a4paper,top=1.5cm,bottom=1.35cm,left=1.6cm,right=1.6cm]{geometry}
\usepackage{microtype}
\usepackage{xcolor}
\usepackage{titlesec}
\usepackage[most]{tcolorbox}
\usepackage{enumitem}
\usepackage{array}
\usepackage{booktabs}
\usepackage{fancyhdr}

\definecolor{navy}{HTML}{1F3A5F}
\definecolor{grey}{HTML}{54606E}
\definecolor{green}{HTML}{2C6E49}
\definecolor{amber}{HTML}{9C6A16}
\definecolor{red}{HTML}{A6322E}
\definecolor{boxgray}{HTML}{F3F5F8}
\definecolor{boxgreen}{HTML}{F1F6F2}
\definecolor{boxedge}{HTML}{D9E0EA}

\setlength{\parindent}{0pt}
\setlength{\parskip}{2.6pt}
\linespread{1.0}

\pagestyle{fancy}\fancyhf{}
\lhead{\small\color{grey}Scella}
\rhead{\small\color{grey}Classement par potentiel -- valeur intrinsèque}
\cfoot{\small\color{grey}\thepage}
\renewcommand{\headrulewidth}{0.3pt}

\titleformat{\section}{\color{navy}\large\bfseries}{}{0pt}{}
\titlespacing{\section}{0pt}{9pt}{3pt}

\newtcolorbox{these}{colback=boxgray,colframe=boxedge,boxrule=0.5pt,arc=2.5pt,
  left=8pt,right=8pt,top=4.5pt,bottom=4.5pt}
\newtcolorbox{cle}{colback=boxgreen,colframe=green!45,boxrule=0.8pt,arc=2.5pt,
  left=8pt,right=8pt,top=5pt,bottom=5pt}

\newcommand{\code}[1]{\texttt{\small #1}}
\newcommand{\best}[1]{{\color{green}\textbf{#1}}}
\newcommand{\sz}[1]{{\small\color{grey}\textsf{#1}}}

\begin{document}

\begin{center}
{\Huge\bfseries\color{navy} Scella}\\[4pt]
{\large\color{grey}\itshape Classement des composants par potentiel -- valeur intrinsèque, hors tout objectif}\\[4pt]
{\small\color{grey} 5 août 2026 \; · \; 4 lentilles intrinsèques \; · \; variante du classement à 5 lentilles}
\end{center}
\vspace{2pt}

\begin{these}
\textbf{Le résultat central, et il est inattendu.} Retirer l'objectif ne change \emph{presque rien} :
sur treize composants, \textbf{un seul échange de place}, entre deux composants à égalité. Le classement
n'était donc pas un artefact du but poursuivi -- il mesurait bien la valeur intrinsèque. Deux choses
bougent néanmoins, et elles comptent : \textbf{le Relais et le reçu deviennent quasiment à égalité}, et
\textbf{le banc de mesure perd son avocat} tout en devenant le composant le plus consensuellement placé
de tout le portefeuille.
\end{these}

\section{1 · Ce qui a été retiré, et pourquoi c'est propre}

Le classement précédent reposait sur cinq lentilles. Quatre d'entre elles -- la \textbf{maturité}
rapportée à l'effort restant, la \textbf{rareté} pure, l'\textbf{adversaire} en revue diligente, et la
\textbf{démontrabilité} devant un ingénieur sceptique -- ne présupposent \emph{aucun} objectif : elles
mesurent ce qu'une chose vaut en elle-même. La cinquième, l'argument d'embauche, était la seule à
dépendre d'un but extérieur. La retirer ne mutile donc pas l'analyse : elle la rend \textbf{neutre quant
à la finalité}. Le classement ci-dessous répond à une question plus simple et plus dure -- \emph{qu'est-ce
qui a le plus de valeur, indépendamment de ce qu'on veut en faire ?}

\section{2 · Le classement intrinsèque}

\renewcommand{\arraystretch}{1.1}
\begin{center}\small
\begin{tabular}{@{}r>{\raggedright\arraybackslash}p{3.2cm}cccc c c c >{\raggedright\arraybackslash}p{5.4cm}@{}}
\toprule
& \textbf{Composant} & \sz{Mat} & \sz{Rar} & \sz{Adv} & \sz{Dém} & \textbf{Moy} & \sz{Éc.} & \sz{$\Delta$} & \textbf{Ce qui le porte, ce qui le plombe} \\
\midrule
\textbf{1} & \textbf{Le Relais} & 2 & 3 & 3 & \best{1} & \textbf{2,25} & 2 & \sz{=} & Original \emph{et} construit \emph{et} servi ; le seul sans caveat \\
\textbf{2} & \textbf{Le reçu, pile A} & \best{1} & 2 & \best{1} & 6 & \textbf{2,50} & 5 & \sz{=} & Scelle le réel ; rend \code{verified:true} sur une chaîne tronquée \\
\textbf{3} & \textbf{Le reçu, pile B} & 5 & \best{1} & 6 & 2 & \textbf{3,50} & 5 & \sz{=} & Vérifieur à dépendances vides \emph{testé} ; rien à vérifier \\
\textbf{4} & \textbf{Le banc de mesure} & 4 & 4 & 4 & 4 & \textbf{4,00} & \best{0} & \sz{=} & Réfute sa propre thèse ; rien n'est rejouable \\
5 & La Couche 2 servie & 3 & 10 & \best{2} & 5 & 5,00 & \textbf{8} & \sz{=} & Seule mesure en argent réel ; catégorie saturée \\
6 & L'essaim $\Phi$ & 7 & 5 & 7 & 3 & 5,50 & 4 & \sz{=} & Tue sa propre borne en source ; concurrence non pure \\
7 & Le Sandwich & 6 & 9 & 5 & 7 & 6,75 & 4 & \sz{=} & Substrat prêt ; rareté seulement dérivée du reçu \\
8 & L'essaim médiateur & 9 & 8 & 9 & 10 & 9,00 & 2 & \sz{=} & Anti-émergence assumée ; catégorie encombrée, 0 ligne \\
9\,= & Fil C, la Colonne & 12 & 7 & 8 & 11 & 9,50 & 5 & \sz{+1} & Objet visé unique ; rareté empruntée, précondition absente \\
9\,= & Le Flexible MLP & 13 & 6 & 10 & 9 & 9,50 & \textbf{7} & \sz{--1} & Seul actif \emph{déductif} ; jamais flashé, hors dépôt \\
11 & Les nano-robots & 8 & 13 & 11 & 8 & 10,00 & 5 & \sz{=} & Réels ; priors \code{sin(i)} auto-qualifiés « non signifiants » \\
12 & Le méta-essaim & 11 & 11 & 12 & 12 & 11,50 & \best{1} & \sz{=} & Volume maximal, valeur réalisée nulle \\
13 & Le runtime nano & 10 & 12 & 13 & 13 & 12,00 & 3 & \sz{=} & Ré-implémente du résolu, et le ré-implémente faux \\
\bottomrule
\end{tabular}
\end{center}

\sz{$\Delta$ = déplacement par rapport au classement incluant la lentille d'embauche. Éc. = amplitude
entre le meilleur et le pire rang.}

\section{3 · Les deux seuls changements réels -- et ce qu'ils disent}

\textbf{Le Relais et le reçu passent à égalité.} C'est le changement de fond. Dans le classement à cinq
lentilles, l'écart entre eux était net (2,00 contre 3,20) ; il tombe ici à \textbf{2,25 contre 2,50},
c'est-à-dire à rien. La raison est instructive : la lentille d'embauche était \emph{celle qui punissait
le plus durement la pile A}, parce qu'un candidat qui vend la vérifiabilité et dont la route de
vérification ne compare jamais la racine se livre un contre-argument. Ce raisonnement est parfaitement
juste pour un entretien -- et parfaitement hors sujet pour évaluer un actif. Retirée cette pénalité de
présentation, \textbf{le reçu qui scelle les missions réellement servies remonte au niveau du Relais}, et
domine même trois lentilles sur quatre : premier en maturité rapportée à l'effort, premier chez
l'adversaire, deuxième en rareté. Il ne perd que sur la démontrabilité, à cause du défaut de troncature
-- lequel est de taille S.

\textbf{Le banc de mesure perd son avocat, et gagne un titre.} Il était deuxième pour l'embauche, sur un
argument fort : savoir réfuter sa propre thèse est un signal rare chez un ingénieur. Hors de ce cadre,
cet argument perd sa force -- personne n'achète un outil interne pour la vertu de son propriétaire. Il
reste donc quatrième, exactement là où il était. Mais il gagne au passage une distinction que personne
d'autre n'a : \textbf{un écart de zéro}. Les quatre lentilles intrinsèques, qui se disputent sur presque
tout, le placent \emph{toutes les quatre} au quatrième rang. C'est le composant le plus consensuellement
situé du portefeuille -- ni surestimé ni sous-estimé par aucune façon de regarder.

\textbf{Et un non-changement qui vaut un résultat.} Fil C et le Flexible MLP échangent leurs places, mais
à égalité stricte (9,50 chacun) : ce n'est pas un mouvement, c'est un départage. \emph{Aucun autre
composant ne bouge d'un rang.} Un classement qui survit intact au retrait de son cadre stratégique est un
classement qui mesurait la chose et non le contexte.

\section{4 · Ce que la lentille adverse voit, et qui ne dépend d'aucun objectif}

Les trois découvertes structurelles du panel sont, elles aussi, indépendantes de toute finalité, et
c'est ce qui fait leur solidité.

\textbf{Le défaut récurrent est l'auto-attestation.} Cinq fois, dans cinq organes indépendants, la boucle
de vérification se referme à l'intérieur du processus qui a produit la chose : le reçu de la pile A
recalcule la racine sans jamais la comparer à la racine scellée ; celui de la pile B n'ancre rien, si
bien que la racine de référence \emph{voyage dans le dossier} tendu à l'auditeur ; le méta-essaim possède
le service qui prouverait à quel moteur on parle, et ce service n'a aucun appelant ; la Couche 2 affiche
« préflight OK » alors que la rédaction des données personnelles est compilée dehors ; le banc gèle ses
constantes mais ne commite aucun script rejouable. Ce ne sont pas cinq chantiers, c'est \textbf{un seul
principe de réparation} : déplacer chaque point de vérification \emph{hors} du producteur.

\textbf{Les quelque 9\,500 tests forment une seule réserve de crédibilité partagée.} Un examinateur qui
attrape \textbf{un seul} faux vert cesse de croire aux 320, aux 245, aux 118 et aux 49 autres, y compris
ceux qui sont irréprochables. Cinq défauts localisés tiennent l'ensemble en otage. Le corollaire inverse
l'ordre des priorités : le meilleur rendement n'est pas de coder le Sandwich, l'essaim médiateur ou la
Colonne -- qui n'ajouteraient que des lignes à défendre -- mais \textbf{d'éteindre les cinq faux verts},
dont trois sont petits et déjà outillés.

\textbf{Le reçu se scinde, et les deux moitiés ne se touchent jamais.} La pile A scelle les missions
réellement servies mais échoue à la première sonde du métier -- tronquer, re-vérifier, obtenir
\code{verified:true} sans détenir la moindre clé. La pile B tient une revendication extraordinaire et
\emph{testée} -- un vérificateur dont l'arborescence de dépendances est littéralement vide, une
recompilation hors ligne de la source vendorée, puis une re-preuve -- mais n'a rien à vérifier et n'ancre
rien. Chacune porte très exactement le défaut que l'autre corrigerait.

\section{5 · Le geste, désormais encore plus évident}

Sans objectif extérieur, la conclusion du panel ne s'affaiblit pas : elle se durcit. Les deux premiers du
classement sont \emph{les deux moitiés d'un même actif}, et le geste qui les réunit est le seul de tout
le dépôt à transformer une addition en multiplication.

\begin{cle}
\textbf{Brancher les reçus de la pile A dans le graphe de la pile B, signer la racine en \emph{ed25519}
-- la bibliothèque est \underline{déjà} dans la crate -- et publier l'ancrage \emph{hors} du dossier.}
Ce qui rend le geste court : les empreintes de contenu de la pile A sont \textbf{déjà sans clé}, donc
n'importe qui peut déjà les recalculer ; seuls le HMAC de maillon et le pli de racine exigent le secret.
\sz{Cryptographie : taille S. Garde et rotation des clés : taille M.}
\end{cle}

L'objet obtenu -- \emph{un reçu de mission multi-modèles re-prouvable par un binaire dont l'auditeur peut
lire la totalité de la source} -- n'existe nulle part ailleurs, et sa valeur ne dépend d'aucun
destinataire particulier. C'est aussi le seul chemin par lequel le profil de souveraineté visé par la
Colonne cesse d'être une idée empruntée.

\textbf{Le classement intrinsèque se lit donc ainsi :} le reçu \emph{réuni} est premier, sans discussion ;
le Relais est premier parmi les composants \emph{tels qu'ils existent}, à égalité avec la moitié servie du
reçu. Tout le reste est loin derrière -- et les quatre derniers (nano-robots, méta-essaim, runtime nano,
et la Colonne tant que sa précondition manque) ne méritent aucun effort avant que les cinq faux verts
soient éteints.

\vspace{3pt}
{\color{grey}\hrule height 0.4pt}
\vspace{2pt}
{\footnotesize\color{grey}Variante du classement à cinq lentilles, la lentille d'embauche retirée. Les
rangs par lentille sont inchangés -- seule la moyenne est recalculée sur les quatre lentilles
intrinsèques. Base factuelle : l'état des composants vérifié dans le code, non dans les documents de
conception.}

\end{document}
